% Variants and Ligatures
\section{Variants and Ligatures}

\subsection{Variant Types}

The Indus script exhibits several types of systematic variation:

\subsubsection{Allographic Variants}
Different physical realizations of the same operational sign:
\begin{itemize}
    \item Simplified vs. detailed forms
    \item Compact vs. extended representations
    \item Site-specific calligraphic styles
\end{itemize}

\subsubsection{Orientational Variants}
Signs that differ primarily in orientation:
\begin{itemize}
    \item Left-facing vs. right-facing animals
    \item Upright vs. inverted symbols
    \item Rotated geometric forms
\end{itemize}

\textbf{Open Question}: Whether orientational variants carry semantic/operational significance or are purely stylistic.

\subsection{Ligatures}

Composite signs formed by combining simpler elements:

\subsubsection{Types of Ligatures}
\begin{enumerate}
    \item \textbf{Additive ligatures}: Two signs placed in proximity without fusion
    \item \textbf{Superposition ligatures}: Signs overlaid or merged
    \item \textbf{Enclosed ligatures}: One sign contained within another
\end{enumerate}

\subsubsection{Functional Interpretation}
Ligatures may represent:
\begin{itemize}
    \item Operational composition (nested functions)
    \item Qualification (attribute assignment)
    \item Space-saving notation
    \item Semantic compounds
\end{itemize}

\subsection{Notation Compression}

Evidence suggests ligatures function as notational shortcuts:
\begin{itemize}
    \item Frequent sign combinations may be written as ligatures
    \item Ligatures typically appear on space-constrained objects (small seals)
    \item Expanded forms occur on tablets and pottery with more available space
\end{itemize}

\subsection{Statistical Treatment}

For distributional analysis, we treat:
\begin{itemize}
    \item Allographic variants as a single operational sign
    \item Orientational variants as distinct signs (pending further evidence)
    \item Ligatures as sign sequences or as atomic signs depending on context
\end{itemize}

This treatment is conservative and biases toward demonstrable distinctions.
